\section{Related Work}

\begin{table}[t]
\centering\small
\setlength{\tabcolsep}{6pt}
\setlength{\aboverulesep}{0pt}
\setlength{\belowrulesep}{0pt}
\setlength{\extrarowheight}{0.4ex}
\begin{tabular}{@{}lccccc@{}}
\toprule
Method & Mimic & Latent & Downstream & Task signal & Real robot \\
\midrule
NPMP~\citep{merel2018neural}                & SVG(0) & BC     & SVG(0)      & RL                & \xmarkr \\
ControlVAE~\citep{yao2022controlvae}        & \multicolumn{2}{>{\columncolor{AnaFill}}c}{MBRL}     & MBRL    & RL                & \xmarkr \\
PhysicsVAE~\citep{won2022physics}           & PPO   & BC$^1$ & PPO         & RL                & \xmarkr \\
PULSE~\citep{luo2024universal}              & PPO   & DAgger & PPO         & RL                & \xmarkr \\
MaskedMimic~\citep{tessler2024maskedmimic}  & A2C   & \multicolumn{2}{>{\columncolor{AnaFill}}c}{DAgger} & Distill     & \xmarkr \\
BeyondMimic~\citep{liao2025beyondmimic}     & PPO   & DAgger & CGD         & Distill         & \cmarkg \\
SONIC~\citep{luo2025sonic}                  & \multicolumn{3}{>{\columncolor{AnaFill}}c}{PPO}      & Distill$^2$ & \cmarkg \\
LATENT~\citep{zhang2026learning}            & PPO   & DAgger & PPO         & RL                & \cmarkg \\
\midrule
\textbf{SuperMimic}          & \multicolumn{3}{>{\columncolor{AnaFill}}c}{MBRL}              & Distill\ /\ RL                & \cmarkg \\
\bottomrule
\end{tabular}
\caption{\textbf{Comparison of various methods and their training strategies across different stages.} The Mimic, Latent, and Downstream columns indicate the training strategy used at each stage. Task signal is what trains the downstream interface: its own task reward (RL) or a reference-tracking signal (Distill). Merged cells denote jointly trained stages, while adjacent unmerged cells indicate sequential training, with each stage trained after freezing the previous one.\protect\newline SVG(0): stochastic value gradients; BC: offline behavior cloning on fixed expert rollouts, as opposed to DAgger's student rollouts; CGD: classifier guidance diffusion; A2C: advantage actor--critic; DAgger: dataset aggregation. $^1$Behavior cloning with a one-step next-state consistency term through a frozen, learned dynamics model. $^2$Goal encoders train jointly with the tracker on the motion-tracking reward and latent-alignment losses; no task-specific reward.}
\label{tab:related}
\end{table}

\textbf{Large-scale physics-based motion tracking.}
Physics-based motion imitation learns controllers that reproduce reference motions under contact-rich rigid-body dynamics.
DeepMimic~\citep{peng2018deepmimic} demonstrated expressive character skills through model-free reinforcement learning, and PHC~\citep{luo2023perpetual} extended this paradigm to large motion sets with robust recovery.
SONIC~\citep{luo2025sonic} further demonstrated that scaling motion data, policy capacity, and training compute produces generalist controllers that transfer to humanoid hardware.
These results establish motion tracking as a practical route to broad whole-body behavior coverage.
However, training these large-scale motion trackers requires extensive simulator rollouts and substantial computation.
SuperMimic addresses this bottleneck by learning a differentiable dynamics model shared across all motions, so policy improvement can exploit analytic gradients rather than model-free rollout estimates.

\textbf{Learned differentiable dynamics for motion control.}
SuperTrack~\citep{fussell2021supertrack} optimizes motion-control policies through a learned differentiable surrogate of the simulator.
ControlVAE~\citep{yao2022controlvae} jointly learns a dynamics model, latent skill representation, and generative controller, while MoConVQ~\citep{yao2024moconvq} extends learned-model training to discrete motion representations over tens of hours of data.
SuperDyno~\citep{lei2025scalable} scales differentiable neural dynamics to tens of thousands of motion clips.
These methods establish the effectiveness of learned dynamics and analytic policy gradients for scalable simulated-character control.
SuperMimic extends this direction to single-GPU training of a large-scale humanoid tracker and to its sim-to-real deployment: its world model conditions on simulator context, while the deployed actor uses noisy proprioceptive observations and a goal input (\S\ref{sec:wm}).
We evaluate how model accuracy, rollout throughput, and differentiable policy updates contribute to training efficiency, and how prediction error behaves as motion coverage grows (\S\ref{sec:scaling}--\ref{sec:wm-properties}), and demonstrate physical deployment of the resulting controller (\S\ref{sec:real-robot}).

\textbf{Model-based policy optimization and robot deployment.}
Dreamer~\citep{hafner2019dream} and DreamerV3~\citep{hafner2025dreamerv3} learn latent dynamics for imagination-based actor--critic learning, while TD-MPC~\citep{hansen2022temporal} combines short-horizon planning with terminal-value prediction.
DayDreamer~\citep{wu2023daydreamer} demonstrates world-model learning directly on physical robots.
SHAC~\citep{xu2022accelerated} uses differentiable simulation for short-horizon actor--critic updates, motivating SuperMimic's critic-bootstrapped objective.
Analytic-gradient robot control can also combine a simplified differentiable surrogate with a high-fidelity forward simulator, as demonstrated for quadruped locomotion~\citep{song2025differentiable}.
SuperMimic learns its surrogate from simulator transitions and backpropagates through short imagined rollouts, retaining compatibility with standard rigid-body simulators.

RWM~\citep{li2025robotic} and RWM-U~\citep{li2025uncertainty} demonstrate that policies trained in learned neural simulators can transfer to physical robots.
Both optimize policies using PPO on imagined rollouts; RWM-U additionally uses uncertainty-penalized offline optimization.
SuperMimic instead optimizes the policy by differentiating the return through its learned dynamics.
RWM and RWM-U pretrain the world model on trajectories from other policies. In contrast, SuperMimic fits its world model during training on data its own policy collects.

\textbf{Reusable motion representations and downstream control.}
Reusable controllers connect motion tracking to downstream goal interfaces (Table~\ref{tab:related}).
NPMP~\citep{merel2018neural} and PhysicsVAE~\citep{won2022physics} learn latent motor representations from expert behavior; PULSE~\citep{luo2024universal} and LATENT~\citep{zhang2026learning} use learned motion representations for downstream reinforcement learning.
MaskedMimic~\citep{tessler2024maskedmimic} supports partial-body and multimodal goals through motion inpainting, while BeyondMimic~\citep{liao2025beyondmimic} uses guided diffusion for goal following.
SONIC~\citep{luo2025sonic} jointly trains goal encoders with its tracker using tracking rewards and latent-alignment losses.
ControlVAE~\citep{yao2022controlvae} optimizes downstream task policies through its learned model while freezing the skill controller.
SuperMimic jointly trains goal-specific encoders and a shared decoder through the same dynamics model (Appendix~\ref{app:unified}).
Goal interfaces can optimize either a reference-tracking reward or their own task reward.
Our preliminary study trains a velocity interface concurrently with the motion-tracking interface, once on each objective (\S\ref{sec:downstream}).
